
\documentclass{article}
\usepackage{colortbl}
\usepackage{makecell}
\usepackage{multirow}
\usepackage{supertabular}

\begin{document}

\newcounter{utterance}

\centering \large Interaction Transcript for game `cladder', experiment `full\_v1.5\_default', episode 916 with Qwen3.5{-}4B.
\vspace{24pt}

{ \footnotesize  \setcounter{utterance}{1}
\setlength{\tabcolsep}{0pt}
\begin{supertabular}{c@{$\;$}|p{.15\linewidth}@{}p{.15\linewidth}p{.15\linewidth}p{.15\linewidth}p{.15\linewidth}p{.15\linewidth}}
    \# & \multicolumn{2}{c}{Player} && \multicolumn{2}{c}{Game Master} \\
    \hline

    \theutterance \stepcounter{utterance}  
    & & & \multicolumn{4}{p{0.6\linewidth}}{
        \cellcolor[rgb]{0.9,0.9,0.9}{
            \makecell[{{p{\linewidth}}}]{
                \texttt{\tiny{[P1$\langle$GM]}}
                \texttt{You are an expert in causal inference. The following question is not a typical commonsense query, but rather a meticulously designed question created by a professor specializing in causal inference, intended to assess the students' mastery of the course content.} \\
\texttt{Imagine a self{-}contained, hypothetical world with only the following conditions, and without any unmentioned factors or causal relationships: Smoking has a direct effect on effort and college admission. Effort has a direct effect on college admission. For nonsmokers and are lazy, the probability of college admission is 47\%. For nonsmokers and are hard{-}working, the probability of college admission is 88\%. For smokers and are lazy, the probability of college admission is 23\%. For smokers and are hard{-}working, the probability of college admission is 54\%. For nonsmokers, the probability of being hard{-}working is 40\%. For smokers, the probability of being hard{-}working is 61\%. If we disregard the mediation effect through effort, would smoking negatively affect college admission?} \\
\texttt{Start your answer with 'Yes' or 'No', followed by additional reasoning or evidence to support your explanation.} \\
            }
        }
    }
     \\ \\

    \theutterance \stepcounter{utterance}  
    & & \multicolumn{4}{p{0.6\linewidth}}{
        \cellcolor[rgb]{0.95,0.95,0.95}{
            \makecell[{{p{\linewidth}}}]{
                \texttt{\tiny{[P1$\leftrightarrow$P]}}
                \texttt{{[}DRAFT{]}} \\
\texttt{Yes} \\
\\ 
\texttt{**Reasoning:**} \\
\\ 
\texttt{To determine whether smoking negatively affects college admission when disregarding the mediation effect through effort, we must calculate the **Total Effect** of smoking on admission probability. In causal inference, the total effect is the difference in the outcome probability between the treated group (smokers) and the control group (nonsmokers), holding the distribution of other variables constant or by comparing the marginal probabilities in the absence of confounding mechanisms that we are instructed to ignore.} \\
\\ 
\texttt{The problem explicitly states to "disregard the mediation effect through effort." This implies we should treat "effort" as a non{-}existent or irrelevant pathway for this specific calculation, effectively comparing the direct causal impact of smoking on admission. We can derive this by comparing the marginal probability of admission for smokers versus nonsmokers based on the provided conditional probabilities.} \\
\\ 
\texttt{1.  **Calculate the marginal probability of admission for Nonsmokers (\$P(\textbackslash{}text\{Admission\} | \textbackslash{}text\{Nonsmoker\})\$):**} \\
\texttt{    *   \$P(\textbackslash{}text\{Admission\} | \textbackslash{}text\{Nonsmoker, Lazy\}) = 0.47\$} \\
\texttt{    *   \$P(\textbackslash{}text\{Admission\} | \textbackslash{}text\{Nonsmoker, Hard{-}working\}) = 0.88\$} \\
\texttt{    *   \$P(\textbackslash{}text\{Hard{-}working\} | \textbackslash{}text\{Nonsmoker\}) = 0.40\$ (Therefore, \$P(\textbackslash{}text\{Lazy\} | \textbackslash{}text\{Nonsmoker\}) = 0.60\$)} \\
\texttt{    *   \$P(\textbackslash{}text\{Admission\} | \textbackslash{}text\{Nonsmoker\}) = (0.47 \textbackslash{}times 0.60) + (0.88 \textbackslash{}times 0.40)\$} \\
\texttt{    *   \$P(\textbackslash{}text\{Admission\} | \textbackslash{}text\{Nonsmoker\}) = 0.282 + 0.352 = 0.634\$ (or 63.4\%)} \\
\\ 
\texttt{2.  **Calculate the marginal probability of admission for Smokers (\$P(\textbackslash{}text\{Admission\} | \textbackslash{}text\{Smoker\})\$):**} \\
\texttt{    *   \$P(\textbackslash{}text\{Admission\} | \textbackslash{}text\{Smoker, Lazy\}) = 0.23\$} \\
\texttt{    *   \$P(\textbackslash{}text\{Admission\} | \textbackslash{}text\{Smoker, Hard{-}working\}) = 0.54\$} \\
\texttt{    *   \$P(\textbackslash{}text\{Hard{-}working\} | \textbackslash{}text\{Smoker\}) = 0.61\$ (Therefore, \$P(\textbackslash{}text\{Lazy\} | \textbackslash{}text\{Smoker\}) = 0.39\$)} \\
\texttt{    *   \$P(\textbackslash{}text\{Admission\} | \textbackslash{}text\{Smoker\}) = (0.23 \textbackslash{}times 0.39) + (0.54 \textbackslash{}times} \\
            }
        }
    }
    & \\ \\

    \theutterance \stepcounter{utterance}  
    & & \multicolumn{4}{p{0.6\linewidth}}{
        \cellcolor[rgb]{0.95,0.95,0.95}{
            \makecell[{{p{\linewidth}}}]{
                \texttt{\tiny{[P1$\leftrightarrow$P]}}
                \texttt{{[}SELF{-}FEEDBACK{]}} \\
\texttt{FORMAT: problem} \\
\texttt{RULES: problem} \\
\texttt{STRATEGY: problem} \\
\texttt{REASON: The answer fails to complete the mathematical calculation, cuts off mid{-}sentence, and does not provide the final "Yes" or "No" conclusion required by the prompt.} \\
            }
        }
    }
    & \\ \\

    \theutterance \stepcounter{utterance}  
    & \multicolumn{4}{p{0.6\linewidth}}{
        \cellcolor[rgb]{0.561,0.737,0.561}{
            \makecell[{{p{\linewidth}}}]{
                \texttt{\tiny{[P1$\rangle$GM]}}
                \texttt{Yes} \\
\\ 
\texttt{**Reasoning:**} \\
\\ 
\texttt{To determine whether smoking negatively affects college admission when disregarding the mediation effect through effort, we must calculate the **Total Effect** of smoking on admission probability. In causal inference, the total effect is the difference in the outcome probability between the treated group (smokers) and the control group (nonsmokers).} \\
\\ 
\texttt{The problem explicitly states to "disregard the mediation effect through effort." This implies we should treat "effort" as a non{-}existent or irrelevant pathway for this specific calculation, effectively comparing the direct causal impact of smoking on admission by comparing the marginal probabilities of admission for smokers versus nonsmokers.} \\
\\ 
\texttt{1.  **Calculate the marginal probability of admission for Nonsmokers (\$P(\textbackslash{}text\{Admission\} | \textbackslash{}text\{Nonsmoker\})\$):**} \\
\texttt{    *   \$P(\textbackslash{}text\{Admission\} | \textbackslash{}text\{Nonsmoker, Lazy\}) = 0.47\$} \\
\texttt{    *   \$P(\textbackslash{}text\{Admission\} | \textbackslash{}text\{Nonsmoker, Hard{-}working\}) = 0.88\$} \\
\texttt{    *   \$P(\textbackslash{}text\{Hard{-}working\} | \textbackslash{}text\{Nonsmoker\}) = 0.40\$ (Therefore, \$P(\textbackslash{}text\{Lazy\} | \textbackslash{}text\{Nonsmoker\}) = 0.60\$)} \\
\texttt{    *   \$P(\textbackslash{}text\{Admission\} | \textbackslash{}text\{Nonsmoker\}) = (0.47 \textbackslash{}times 0.60) + (0.88 \textbackslash{}times 0.40)\$} \\
\texttt{    *   \$P(\textbackslash{}text\{Admission\} | \textbackslash{}text\{Nonsmoker\}) = 0.282 + 0.352 = 0.634\$ (or 63.4\%)} \\
\\ 
\texttt{2.  **Calculate the marginal probability of admission for Smokers (\$P(\textbackslash{}text\{Admission\} | \textbackslash{}text\{Smoker\})\$):**} \\
\texttt{    *   \$P(\textbackslash{}text\{Admission\} | \textbackslash{}text\{Smoker, Lazy\}) = 0.23\$} \\
\texttt{    *   \$P(\textbackslash{}text\{Admission\} | \textbackslash{}text\{Smoker, Hard{-}working\}) = 0.54\$} \\
\texttt{    *   \$P(\textbackslash{}text\{Hard{-}working\} | \textbackslash{}text\{Smoker\}) = 0.61\$ (Therefore, \$P(\textbackslash{}text\{Lazy\} | \textbackslash{}text\{Smoker\}) = 0.39\$)} \\
\texttt{    *   \$P(\textbackslash{}text\{Admission\} | \textbackslash{}text\{Smoker\}) = (0.23 \textbackslash{}times 0.39) + (0.54 \textbackslash{}times 0.61)\$} \\
\texttt{    *   \$P(\textbackslash{}text\{Admission\} | \textbackslash{}text\{Smoker\}) = 0.0897 + 0.3} \\
            }
        }
    }
    & & \\ \\

    \theutterance \stepcounter{utterance}  
    & & & \multicolumn{2}{p{0.3\linewidth}}{
        \cellcolor[rgb]{0.9,0.9,0.9}{
            \makecell[{{p{\linewidth}}}]{
                \texttt{\tiny{[GM$|$GM]}}
                \texttt{Yes} \\
            }
        }
    }
    & & \\ \\

    \theutterance \stepcounter{utterance}  
    & & & \multicolumn{2}{p{0.3\linewidth}}{
        \cellcolor[rgb]{0.9,0.9,0.9}{
            \makecell[{{p{\linewidth}}}]{
                \texttt{\tiny{[GM$|$GM]}}
                \texttt{yes} \\
            }
        }
    }
    & & \\ \\

    \theutterance \stepcounter{utterance}  
    & & & \multicolumn{2}{p{0.3\linewidth}}{
        \cellcolor[rgb]{0.9,0.9,0.9}{
            \makecell[{{p{\linewidth}}}]{
                \texttt{\tiny{[GM$|$GM]}}
                \texttt{game\_result = WIN} \\
            }
        }
    }
    & & \\ \\

\end{supertabular}
}

\end{document}
